% ==============================================================
\section{Conclusion}\label{sec:conclusion}
% ==============================================================

This paper proposed \textbf{Multi-Dimensional Apriori-Window},
which extends Apriori-Window~\cite{aprioriwindow}
to arbitrary $d$-dimensional grids.
Through three techniques---one-dimensional multiset reduction
per fixed spatial start point,
duplicate-aware $t$-axis stride,
and singleton interval map intersection pruning---the method
exactly enumerates all dense patterns
together with their arbitrarily-shaped dense regions.

Experiments demonstrated that:
(1)~on synthetic data with four embedding shapes,
the proposed method outputs the exact solution,
while existing methods achieved at most
F1\,$=0.917$ and Jaccard\,$=0.190$;
(2)~linear scaling was confirmed
for $N = 10^3$--$10^8$,
processing $10^8$ transactions within 10 minutes;
(3)~on $d{=}3$--$4$ real-world data from crime, complaints,
and earthquake domains,
the method yielded interpretable insights
consistent with domain knowledge,
such as compound crime hotspots
and depth-stratified seismic clusters.

\paragraph{Limitations.}
(1)~The method assumes grid discretization,
and the choice of bin width and $\bm{W}$ affects results.
(2)~As $W_{\mathrm{sp}}$ or $d$ increases,
the $O(N \cdot W_{\mathrm{sp}})$ cost of multiset construction grows,
requiring dimensionality reduction or coarse binning
in high-dimensional settings.
(3)~The method is specialized for detection and description;
causal explanation and association with external events
are outside its scope.

\paragraph{Future work.}
(1)~\emph{Adaptive parameter selection}:
A framework that automatically adjusts $\bm{W}$ and $\theta$
based on local data statistics
(density distributions, spatial autocorrelation)
is practically important.
(2)~\emph{Further acceleration}:
Construction of temporal multisets per spatial start point
is the main computational bottleneck.
Prefix-sum-based $O(1)$ window aggregation
and incremental update along spatial axes (slice differencing)
to eliminate redundant computation are promising.
